\mode*
\section{Method}

This study has two parts, mirroring the companion paper
\autocite{VTProgMisconceptions}.
The first is a \emph{literature review} that maps learners' misconceptions
and difficulties with LaTeX, answering \cref{rq:misconceptions}:

\rqmisconceptions*

The second is an \emph{analysis} of those difficulties through the lens of
variation theory: they are treated as phenomenographic observations from
which we identify critical aspects (\cref{rq:aspects}) and construct
patterns of variation (\cref{rq:patterns}):

\rqaspects*
\rqpatterns*

\Cref{rq:misconceptions} is answered by the review (and, given the
literature gap, by our own course data); \cref{rq:aspects,rq:patterns} are
answered analytically, and their answers are validated by argument
(inference to the best explanation), not yet empirically.

\subsection{Literature review}
\label{sec:method-litreview}

The review has two phases.
The first phase is a seed round: a claim-driven search conducted for the
course-development work this paper grows out of, with every retained
reference verified against at least its abstract and recorded with full
provenance; the protocol and queries are in \cref{app:protocol}.
That round found \emph{no} studies of LaTeX-specific misconceptions---only
the efficiency comparison \autocite{Knauff2014Latex}.
The second phase is a systematic, tool-supported multi-database search
with the \texttt{scholar} tool, as in the companion papers
\autocite{VTProgMisconceptions,VTDebug}; given the phase-one gap it also
covered adjacent literatures---markup languages in education, compiler
error messages, and end-user programming of documents.
This round confirmed the gap (no LaTeX-specific learner study) and
retained the transferable adjacent-literature sources on which
\cref{sec:errors,sec:markup} draw; two further rounds extended it to
field-tagged queries in Web of Science and Scopus.
A final round asked two prior questions: which aspects of document
preparation the adjacent literatures document as difficult---so that each
candidate aspect of \cref{sec:compilation,sec:errors,sec:markup} can be
marked as read off the literature, off our course, or posited---and
whether learners' understanding of LaTeX or markup has been measured at
all, which \cref{sec:method-instruments} relies on.
All rounds are documented in full in \cref{app:protocol}.

\subsection{Course data}
\label{sec:method-course-data}

% TODO(#2): Because the literature is thin, plan the course-data collection
% explicitly: datintro's LaTeX lab grading output and tutoring
% observations as sources of documented difficulties, collected and coded
% under an agreed scheme. Keep the classification sources separate from
% any data they are later correlated with.

\subsection{Variation-theory analysis}
\label{sec:method-vt-analysis}

For each documented difficulty we ask: which aspect had the learner not
discerned?
Aspects recurring across difficulties are candidate critical aspects.
For each critical aspect we construct patterns of variation---contrast,
generalisation, fusion---following the companion paper
\autocite{VTProgMisconceptions}.
LaTeX affords the vary-and-observe loop
\autocite{Svensson2020Affordances,ThuneEckerdal2009Variation,
ThuneEckerdal2018Lab} in a particularly clean form: change one thing in
the source, recompile, observe the output---source diff against output
diff.

\subsection{Course instruments}
\label{sec:method-instruments}

To make the analysis empirically testable in a future study, we prepare a
pre/post-test that measures whether students discern the candidate
critical aspects of \cref{sec:compilation,sec:errors,sec:markup} before
and after the LaTeX module.
Two short quizzes---administered at the start and the end of the
module---share the same items: one graded choice item per candidate
critical aspect, each distractor expressing the behaviour of a learner who
has not discerned that aspect.
The start-of-module scores measure what students bring, and the gain to
the end-of-module scores measures what the module added; following
\citeauthor{NCOL}, the stems do not point out the aspect under test
\autocite[p.~91]{NCOL}.
Participation is opt-in and pseudonymous: the start quiz opens with a
consent item, and the instrument is set up in the learning platform, and
its results prepared for analysis, by the literate program of
\cref{app:quiz}.
This complements the course-data collection
(\cref{sec:method-course-data}): the course data documents the
difficulties that ground the aspects, while the pre/post-test would later
test whether teaching against them makes the aspects discernible.

The pre/post-test complements the course data of
\cref{sec:method-course-data} as the second of two sources, and with the
literature gap confirmed (\cref{app:protocol}) the course data is
promoted from supplement to primary evidence: the LaTeX lab's grading
already checks semantic-markup use, so each such check failure is a
documented instance of the presentational-markup misconception
(\cref{sec:markup}), and the submitted sources make hand-formatted
headings, manual numbering and unresolved errors statically detectable.
The two sources are kept methodologically separate to avoid
circularity: misconception codes derived from the assessment traces are
never correlated with those same traces; the pre/post-test is the
independent measure of discernment.
Participation is opt-in through the consent item of the pre-test, which
covers course submissions; grading runs for all students regardless,
and research extraction filters by consent afterwards.
